\documentclass[aspectratio=169]{beamer}
\usepackage[style=ieee,backend=biber]{biblatex}
\addbibresource{reference.bib}
\usepackage{colortbl,tabularx,mathrsfs,calligra}
\usepackage{amsmath,amsfonts,amssymb,amsthm}
\usepackage{ragged2e}
\usepackage[bahasa]{babel}
\usepackage{tikz}
\usepackage{caption}
\usepackage{wrapfig}
\usepackage{multirow}
\usepackage{multicol}
\usepackage{array}
\usepackage{pgfplots, tkz-euclide,calc}
\pgfplotsset{compat=1.18}
\usepackage{listings}

\graphicspath{{C:/Users/teoso/OneDrive/Documents/Tugas Kuliah/Template Math Depart/}{./foto/}}

\definecolor{HIMAmuda}{HTML}{01D1FD}
\definecolor{HIMAtua}{HTML}{02016A}
\definecolor{HIMAabu}{HTML}{CBCBCC}
\definecolor{PastelGreen}{HTML}{77DD77}

\usetheme{Madrid}

\setbeamercolor{palette primary}{bg=HIMAtua,fg=white}
\setbeamercolor{palette secondary}{bg=HIMAmuda,fg=black}
\setbeamercolor{palette tertiary}{bg=HIMAabu,fg=black}
\setbeamercolor{palette quaternary}{bg=HIMAmuda,fg=white}
\setbeamercolor{structure}{fg=HIMAmuda} % itemize, enumerate, etc
\setbeamercolor{section in toc}{fg=HIMAtua} % TOC sections
\setbeamercolor{bibliography item}{parent=palette secondary}
\setbeamercolor*{bibliography entry author}{parent=section in toc}

\usetikzlibrary{shapes.geometric, arrows}

\tikzstyle{startstop} = [ellipse, minimum width=1cm, minimum height=1cm,text centered, draw=black, fill=red!30]
\tikzstyle{process} = [rectangle, minimum width=2cm, minimum height=1cm, text centered, draw=black, fill=blue!30]
\tikzstyle{decision} = [diamond, minimum width=1cm, minimum height=1cm, text centered, draw=black, fill=blue!50]
\tikzstyle{arrow} = [thick,->,>=stealth]

\newcolumntype{L}[1]{>{\raggedright\let\newline\\\arraybackslash\hspace{0pt}}m{#1}}
\newcolumntype{C}[1]{>{\centering\let\newline\\\arraybackslash\hspace{0pt}}m{#1}}
\newcolumntype{R}[1]{>{\raggedleft\let\newline\\\arraybackslash\hspace{0pt}}m{#1}}

\usefonttheme{professionalfonts}
\setbeamertemplate{theorems}[numbered]
\setbeamertemplate{bibliography item}{\insertbiblabel}
% \setbeamercovered{transparent}


\theoremstyle{definition}
% \numberwithin{subsection}{section}
\newtheorem{definisi}{Definisi}
\newtheorem{teorema}{TEOREMA}
\newtheorem{contoh}{Contoh}
\newtheorem{latihan}{Latihan}
\newcommand{\R}{\mathbb{R}}
\newcommand{\N}{\mathbb{N}}
\newcommand{\Z}{\mathbb{Z}}
\newcommand{\C}{\mathbb{C}}


\AtBeginEnvironment{contoh}{%
\setbeamercolor{block title}{use=example text,fg=white,bg=example text.fg!75!black}
\setbeamercolor{block body}{parent=normal text,use=block title example,bg=block title example.bg!10!bg}
}
\AtBeginEnvironment{definisi}{
\setbeamercolor{block title}{fg=white,bg=HIMAtua}
\setbeamercolor{block body}{parent=normal text,bg=HIMAtua!30!white}
}
\AtBeginEnvironment{latihan}{%
  \setbeamercolor{block title}{fg=white,bg=yellow!50!orange} % Set title background to pastel green and text to white
  \setbeamercolor{block body}{parent=normal text,bg=yellow!30!white} % Set body background to a lighter pastel green
}

\date{Sabtu, 25 April 2026}
\title[OSK Matematika]{Pembinaan OSN-K Matematika}
\author[Tew \& Jun]{Teosofi Hidayah Agung\\Junika Irdia Indi Astudin}
\institute{MAN 1 Pasuruan}

\begin{document}
\begin{frame}
  \titlepage
\end{frame}

% \begin{frame}{Daftar Isi}
%   \tableofcontents
% \end{frame}

\section{Panduan OSN}
\begin{frame}{\insertsection}
  OSN-K Matematika akan diadakan pada tanggal 18 Juni 2026. Soal yang akan diujikan terdiri dari 10 soal pilihan ganda dan 10 soal isian singkat. Peserta akan diberikan waktu 2 jam 30 menit untuk menyelesaikan semua soal.

  \vspace{0.5cm}

  Panduan Lengkap OSN dapat diakses melalui tautan berikut:
  \begin{center}
    \color{blue}\href{https://pusatprestasinasional.kemendikdasmen.go.id/event/riset-dan-inovasi/sma/olimpiade-sains-nasional-2026-2026-sma}{pusatprestasinasional.kemendikdasmen.go.id}
  \end{center}
\end{frame}

\section{Silabus}
\subsection{Aljabar \& Geometri}
\begin{frame}{\insertsection}{\insertsubsection}
  \begin{columns}[t]
    \begin{column}{0.5\textwidth}
      \begin{table}[H]
        \centering
        \begin{tabular}{|>{\columncolor{HIMAtua!20}}c|p{6cm}|}
          \hline
          \rowcolor{HIMAmuda!50}
          \textbf{No} & \textbf{Aljabar}                    \\ \hline
          A1          & Identitas aljabar                   \\ \hline
          A2          & Barisan aritmatika/geometri         \\ \hline
          A3          & Ketaksamaan (AM-GM, Cauchy-Schwarz) \\ \hline
          A4          & Polinom                             \\ \hline
          A5          & Akar \& algoritma pembagian         \\ \hline
        \end{tabular}
      \end{table}
    \end{column}
    \begin{column}{0.5\textwidth}
      \begin{table}[H]
        \centering
        \begin{tabular}{|>{\columncolor{HIMAtua!20}}c|p{6cm}|}
          \hline
          \rowcolor{HIMAmuda!50}
          \textbf{No} & \textbf{Geometri}        \\ \hline
          G1          & Angle chasing            \\ \hline
          G2          & Garis berat/sumbu/tinggi \\ \hline
          G3          & Lingkaran luar/dalam     \\ \hline
          G4          & Segi empat tali busur    \\ \hline
          G5          & Titik kuasa              \\ \hline
        \end{tabular}
      \end{table}
    \end{column}
  \end{columns}
\end{frame}

\subsection{Teori Bilangan \& Kombinatorika}
\begin{frame}{\insertsection}{\insertsubsection}
  \begin{columns}[t]
    \begin{column}{0.5\textwidth}
      \begin{table}[H]
        \centering
        \begin{tabular}{|>{\columncolor{HIMAtua!20}}c|p{6cm}|}
          \hline
          \rowcolor{HIMAmuda!50}
          \textbf{No} & \textbf{Kombinatorika}                \\ \hline
          K1          & Basic counting (permutasi, kombinasi) \\ \hline
          K2          & Koefisien binomial                    \\ \hline
          K3          & Peluang                               \\ \hline
          K4          & Prinsip bijeksi                       \\ \hline
          K5          & Teori himpunan                        \\ \hline
        \end{tabular}
      \end{table}
    \end{column}
    \begin{column}{0.5\textwidth}
      \begin{table}[H]
        \centering
        \begin{tabular}{|>{\columncolor{HIMAtua!20}}c|p{6cm}|}
          \hline
          \rowcolor{HIMAmuda!50}
          \textbf{No} & \textbf{Teori Bilangan} \\ \hline
          T1          & Keterbagian             \\ \hline
          T2          & FPB/KPK                 \\ \hline
          T3          & Bilangan prima          \\ \hline
          T4          & Persamaan Diophantine   \\ \hline
          T5          & Kongruensi modulo       \\ \hline
          T6          & Teorema Fermat \& Euler \\ \hline
        \end{tabular}
      \end{table}
    \end{column}
  \end{columns}
\end{frame}


\section{\textit{Plan} Pembinaan}
\begin{frame}{\insertsection}
  Pembinaan OSN-K Matematika akan dilaksanakan selama 8 minggu, setiap hari Sabtu jam 09:30-11:30.
  \begin{table}[H]
    \begin{tabular}{|>{\columncolor{HIMAtua!20}}c|c|c|c|l|c|}
      \hline
      \rowcolor{HIMAmuda!50}
      \textbf{No} & \textbf{Tanggal} & \textbf{Pelaksanaan} & \textbf{Pemateri} & \textbf{Materi} & \textbf{SubMateri} \\ \hline
      1           & 25 April         & Offline              & Junika            & Kombinatorika   & -                  \\ \hline
      2           & 2 Mei            & Offline              & Teosofi           & Aljabar         & -                  \\ \hline
      3           & 9 Mei            & Online               & Junika            & Teori Bilangan  & -                  \\ \hline
      4           & 16 Mei           & Online               & Teosofi           & Geometri        & -                  \\ \hline
      5           & 23 Mei           & Online               & Junika            & Kombinatorika   & -                  \\ \hline
      6           & 30 Mei           & Online               & Teosofi           & Aljabar         & -                  \\ \hline
      7           & 7 Juni           & Offline              & Junika            & Teori Bilangan  & -                  \\ \hline
      8           & 14 Juni          & Offline              & Teosofi           & Geometri        & -                  \\ \hline
    \end{tabular}
  \end{table}

  \textit{Catatan: Jadwal dapat berubah}
\end{frame}

\end{document}
